\documentclass[11pt, leqno]{article}

\usepackage{amsmath, amsthm, amssymb, tikz, upgreek,
    mathrsfs, hhline, mathtools, biblatex}

\addbibresource{bibliography.bib} 

\newtheorem{theorem}{Theorem}[section]
\newtheorem{lemma}[theorem]{Lemma}
\newtheorem{proposition}[theorem]{Proposition}
\newtheorem{corollary}[theorem]{Corollary}
\newtheorem{definition}[theorem]{Definition}
\newtheorem{conjecture}[theorem]{Conjecture}

\newcommand*\circled[1]{\tikz[baseline=(char.base)]{
            \node[shape=circle,draw,inner sep=2pt] (char) {#1};}}

            \newcommand\numberthis{\addtocounter{equation}{1}\tag{\theequation}}

\makeatletter
\newcommand*\bigcdot{\mathpalette\bigcdot@{1}}
\newcommand*\bigcdot@[2]{\mathbin{\vcenter{\hbox{\scalebox{#2}{$\m@th#1\bullet$}}}}}
\makeatother

\DeclarePairedDelimiter\floor{\lfloor}{\rfloor}

% \newcommand{\vk}{\varkappa}
\newcommand{\Li}{\mathrm{Li}}
\newcommand{\Di}{\mathcal{D}}
\newcommand{\Res}{\text{Res}}

\begin{document}

\title{An Intuitive Introduction to the First Hardy-Littlewood Conjecture}
\author{Logan McBroom}
\date{}
\maketitle

\begin{abstract}
    We encode relative primality information with respect to k-tuples in
    extensions of common arithmetic functions. 
    This is used to motivate the residue counting function of
    the first Hardy-Littlewood conjecture, and leads
    to the singular series.
    We also discuss the singular series as a probabilistic adjustment.
\end{abstract}

\section{K-Tuples}

\begin{definition}
Define $\mathbb{K}$, the set of \textbf{k-tuples}, by
$$\mathbb{K} := \mathscr{P}(\mathbb{Z}).$$
\end{definition}

This is a non-standard definition.
As the name indicates, k-tuples are usually defined as tuples,
that is, finite sequences.
If we require finiteness, we will manually indicate it,
and the sequence order of k-tuples is never needed.
There is also the indication that `k-tuple' should mean there are $k$
elements in the set, and this is usually the case.
This is convenient, but also confusing when juggling many k-tuples of
multiple sizes, so we take the name k-tuple to simply be a set
of integers, and describe the size of a k-tuple $A$ (or any other set) 
by $\#A$.

Our goal is to view traditional number theoretic concepts
as being about the k-tuple $\{0\}$, and to attempt to extend those results
in a natural way to all k-tuples.

\begin{definition}
For $A,B \in\mathbb{K}$, we define
$$A+B := \{a+b : a\in A, b\in B \}.$$
\end{definition}

For $j\in\mathbb{Z}$ and $A\in\mathbb{K}$, define for convenience 
$j+A := \{j\} + A$.
That is, $j+A$ is the translation of $A$ by $j$.

We are concerned with congruence between elements of k-tuples,
but the usual definition of congruence is insufficient, as it describes only
exact incidence between numbers relative to a modulus.
We are instead concerned with a variety of types of incidence between
translated k-tuples and the integers relatively prime to some modulus.
The following definition handles the ``modulus'' aspect of that need.

\begin{definition}
Given a k-tuple $A$ and a natural number $n$, define
$$A_n := \{a \bmod n : a\in A\}$$
\end{definition}

% Note that if $\vk(n)=C$ for some $n$, $C$ is periodic with period $n$, though again, it may possess a smaller period than $n$
% as will be seen in Proposition \ref{vk primes}.
% For a k-tuple $A$, we can then speak of each element of $A$ being coprime to some $n$ by saying $A\subset\vk(n)$.

% This definition has the issue that for some A, (A,n) is neither 1 nor not 1, which is logically confusing.
% \begin{definition}
%     Let $A$ be a k-tuple and let $n$ be a natural number.
%     If for all $a\in A$ we have $(a,n)=1$, we say $A$ and $n$ are coprime, and write $(A,n)=1$.

%     If for all $a\in A$ we have $(a,n)\not=1$, we write $(A,n)\not= 1$.
% \end{definition}

% \begin{proposition}\label{vk mult}
% For any natural numbers $m$ and $n$
%     $$\vk_{mn}=\vk_m \cap \vk_n.$$
% \end{proposition}
% \begin{proof}
%     For $m,n\in\mathbb{N}$, we have
%     \begin{align*}
%         \vk_{mn} 
%         &= \{ k\in\mathbb{Z}_{mn} : (k,mn)=1 \} \\
%         &= \{ k\in\mathbb{Z} : (k,m)=1 \text{ and } (k,n)= 1 \} \\
%         &= \{ k\in\mathbb{Z} : (k,m)=1 \}
%             \cap \{ k\in\mathbb{Z} : (k,n)= 1 \} \\
%         &= \vk(m) \cap \vk(n).
%     \end{align*}
% \end{proof}
% \begin{corollary}\label{vk subset split}
%     For any natural numbers $m$ and $n$, and any k-tuple $A$, we have
%     $$A \subset \vk(mn)  \Leftrightarrow A\subset \vk(m) \textnormal{ and } A \subset \vk(n).$$
% \end{corollary}
% \begin{proof}
%     Let $m,n\in\mathbb{N}$ and let $A\in\mathbb{K}$.
%     Then by the proposition,
%     \begin{align*}
%         &A\subset \vk(mn) \\
%         \Leftrightarrow &A\subset (\vk(m)\cap\vk(n))\\
%         \Leftrightarrow &A\subset \vk(m) \text{ and } A\subset\vk(n).
%     \end{align*}
% \end{proof}

% \begin{proposition}\label{vk primes}
%     For any prime $p$ and natural number $k$, we have
%     $$\vk(p^k) = \vk(p).$$
% \end{proposition}
% \begin{proof}
%     This is immediate from $(n,p^k)=1 \Leftrightarrow (n,p)=1$.
% \end{proof}

% I don't think this is needed
% \begin{proposition}
%     For any natural number $n$ with prime factorization 
%     $$n=\prod_{p_i|n}p_i^{a_i},$$
%     we have
%     $$\vk(n) = \vk\left(\prod_{p_i|n}p_i\right).$$
% \end{proposition}
% \begin{proof}
%     We have by propositions \ref{vk mult} and \ref{vk primes},
%     \begin{align*}
%         \vk(n) 
%         &= \vk\left(\prod_{p_i|n}p_i^{a_i}\right) \\
%         &= \bigcap_{p_i|n}\vk(p_i^{a_i}) \\
%         &= \bigcap_{p_i|n}\vk(p_i) \\
%         &= \vk\left(\prod_{p_i|n}p_i\right). \\
%     \end{align*}
% \end{proof}

% =======================================================================================================================================
\section{The Totient Function}

We now extend Euler's totient function $\varphi$ to k-tuples.
Instead of counting which natural numbers up to $n$ are relatively 
prime to $n$, we count the translates of a k-tuple $A$ by $j\in\mathbb{Z}_n$ 
for which all elements are coprime to $n$.
We may also consider counting many other types of occurences, and will do so
in a later section.
We graph the counting process below for $A=\{0,2\}$ and $n=6$.
The top grid represents a translate of $A$, and the bottom
are those integers up to 6 coprime to 6.
Recalling that those coprime integers are periodic with period 6,
we simply wrap any elements of the translated $A$ which fall
off the end back to 0.

\begin{enumerate}
    \item[0] \begin{center} \begin{tabular}{ |c|c|c|c|c|c| } 
            \hline
            $\bigcdot$ & $\ $ & $\bigcdot$ & $ $ & $ $ & $\ $ \\
            \hline
            $ $ & $\bigcdot$ & $ $ & $\ $ & $\ $ & $\bigcdot$ \\
            \hline
        \end{tabular} \end{center}
    \item[1] \begin{center} \begin{tabular}{ |c|c|c|c|c|c| } 
        \hline
        $ $ & $\bigcdot$ & $ $ & $\bigcdot$ & $ $ & $\ $ \\
        \hline
        $\ $ & $\bigcdot$ & $\ $ & $\ $ & $\ $ & $\bigcdot$ \\
        \hline
    \end{tabular} \end{center}
    \item[2] \begin{center} \begin{tabular}{ |c|c|c|c|c|c| } 
        \hline
        $ $ & $\ $ & $\bigcdot$ & $ $ & $\bigcdot$ & $\ $ \\
        \hline
        $\ $ & $\bigcdot$ & $ $ & $\ $ & $\ $ & $\bigcdot$ \\
        \hline
    \end{tabular}\\ \end{center}
    \item[3] \begin{center} \begin{tabular}{ |c|c|c|c|c|c| } 
        \hline
        $ $ & $\ $ & $ $ & $\bigcdot$ & $ $ & $\bigcdot$ \\
        \hline
        $\ $ & $\bigcdot$ & $\ $ & $\ $ & $\ $ & $\bigcdot$ \\
        \hline
    \end{tabular} \end{center}
    \item[4] \begin{center} \begin{tabular}{ |c|c|c|c|c|c| } 
        \hline
        $\bigcdot$ & $ $ & $ $ & $ $ & $\bigcdot$ & $\ $\\
        \hline
        $ $ & $\bigcdot$ & $\ $ & $\ $ & $\ $ & $\bigcdot$\\
        \hline
    \end{tabular} \end{center}
    \item[5] \begin{center} \begin{tabular}{ |c|c|c|c|c|c| } 
        \hline
        $ $ & $\bigcdot $ & $ $ & $ $ & $ $ & $\bigcdot$\\
        \hline
        $\ $ & $\bigcdot$ & $\ $ & $\ $ & $\ $ & $\bigcdot$\\
        \hline
    \end{tabular} \end{center}
\end{enumerate}

There is only a single occurrence, on table 5, of complete containment,
so our totient should return 1 for this $A$ and $n$.

\begin{definition}[K-tuple totient]~ \\
    For a k-tuple $A$, define $\varphi_A$ by:
    \begin{align*}
        \varphi_A(1) &:= 1 \\
        \varphi_A(n) &:= \#\{j\in\mathbb{Z}_n : 
            \forall\ a \in A,\ (j+a,n)=1 \}
    \end{align*}
\end{definition}

\begin{theorem}\label{phi multiplicative}
    For any k-tuple $A$, $\varphi_A$ is multiplicative.
\end{theorem}
\begin{proof}
    Throughout, for any integer $j$, let $j_n$ be the class represented 
    by $j$ in $\mathbb{Z}_n$.
    Let $A$ be a k-tuple, 
    and let $m$ and $n$ be coprime natural numbers.
    For each natural number $b$, define the set $P(b)$ by
    $$P(b) = \{j_b\in\mathbb{Z}_b : \forall\ a \in A,\ (j_b+a,b) = 1 \}.$$
    Then by definition $\#P(a)=\varphi_A(a)$. 

    Now let $\Gamma: P(mn) \longmapsto P(m)\times P(n)$ be given by
    $$\Gamma(j_{mn}) = (j_m, j_n).$$
    We will show each of the following:

    \begin{enumerate}
        \item \textbf{$\Gamma$ is well defined.}
            Let $j_{mn}\in P(mn)$.
            Then $\forall\ a\in A,\ (j+a,mn) = 1$, so
            $\forall\ a\in A,\ (j+a,m) = 1$.
            Proceeding likewise for $n$, $(j_m, j_n) \in P(m)\times P(n)$,
            so $\Gamma$ is well defined.

        \item \textbf{$\Gamma$ is injective.}
            Let $a_{mn},b_{mn}\in P(mn)$ such that $\Gamma(a_{mn})=\Gamma(b_{mn})$.
            Then $a_m = b_m$ and $a_n = b_n$. 
            Because $m$ and $n$ are coprime, $a_{mn} = b_{mn}$.
            Thus, $\Gamma$ is injective.
            
        \item \textbf{$\Gamma$ is surjective.}
            Let $(b_m,c_n)\in P(m)\times P(n)$.
            By the Chinese Remainder Theorem, there exists $d_{mn}\in\mathbb{Z}_{mn}$
            such that $d_m=b_m$ and $d_n=c_n$.
            It then suffices to show $d_{mn}\in P(mn)$, for this gives
            $\Gamma(d_{mn}) = (b_m,c_n)$.

            Let $a\in A$. 
            We have
            \begin{align*}
                & b_m\in P(m) \\
                \implies& (a+b_m,m) = 1 \\
                \implies& (a+d_m,m) = 1 \\
                \implies& (a+d_{mn},m) = 1.
            \end{align*}
            
            Proceeding likewise for $c_n$, we have $(a+d_{mn},n) = 1$.
            Combining these gives $(a+d_{mn},mn) = 1$, so 
            $d_{mn}\in P(mn)$ and $\Gamma$ is surjective.
    \end{enumerate}

    Thus, $\Gamma$ is a bijection, so
    \begin{align*}
    \varphi_A(mn)
    &=\#P(mn) \\
    &=\#(P(m)\times P(n)) \\
    &=\#P(m) \cdot \#P(n) \\
    &=\varphi_A(m)\cdot\varphi_A(n).
    \end{align*}
\end{proof}

% ================================================================================================================================
\section{The Unit Function}
The singular unit function $u(n)$ is defined by $u(n)=1$ for all $n$.
It is surprising that such a trivial function has a non-trivial extension to k-tuples.
\begin{definition}
    For any k-tuple $A$ and prime $p$, define $\nu_A(p)$ by:
    $$\nu_A(p) = \#A_p$$
    We extend $\nu_A$ to all natural numbers by complete multiplicativity.
\end{definition}

We have immediately from complete multiplicativity that $\nu_A(1)=1$.
The definition for primes tells us that $\nu_A$ gives the number of residues 
the elements of $A$ fill modulo $p$.
It's usage on primes in the following theorem to describe extensions of $\varphi$ 
is known, at least in some well studied cases.

\begin{theorem}\label{phi prime powers}
    For any k-tuple $A$, prime $p$, and natural number $k$ we have:
    $$\varphi_A( p^k ) = p^k - \nu_A(p)p^{k-1}.$$
\end{theorem}
\begin{proof}
    By definition we have
    \begin{align*}
        \varphi_A(p^k) 
        &= \#\{j \in \mathbb{Z}_{p^k} : \forall\ a \in A,\ (j+a,p^k)=1 \} \\
        &= \#\{j \in \mathbb{Z}_{p^k} : \forall\ a \in A_{p^k},\ (j+a,p^k)=1 \}. \\
    \end{align*}
    If this condition is satisfied by some $j$, it will be satisfied
    by $j+pm$ for any $m$, and there are $p^{k-1}$ multiples of $p$ in $\mathbb{Z}_{p^k}$.
    In other words, as $A_{p^k}$ is translated through $\mathbb{Z}_{p^k}$,
    we gain the same amount towards our total count as we pass through
    each period of size $p$, and there are $p^{k-1}$ such periods.
    We then have
    \begin{align*}
        \varphi_A(p^k) 
        &= p^{k-1}\#\{j \in \mathbb{Z}_p : \forall\ a \in A_p,\ (j+a,p)=1 \} \\
        &= p^{k-1}(p-\#\{j \in \mathbb{Z}_p : \exists\ a \in A_p,\ (j+a,p)\not=1 \}) \\
    \end{align*}
    Now, $(j+a,p)\not=1$ exactly once per $a\in A_p$ as $j$ ranges over $\mathbb{Z}_p$
    (when $j+a \equiv 0 \bmod p$).
    Thus, the size of the set in the previous equation is $\#A_p$, which
    gives the required formula.
\end{proof}
% \begin{proof}[old proof for old defs]
%     By definition
%     \begin{align*}    
%         \nu_A(p)&=\#\{ m\in |A|_p : 0\leq m < p \} \\
%         \Rightarrow p-\nu_A(p)&=\#\{ m\in\mathbb{Z} : 0\leq m < p \textnormal{ and } m \not\in |A|_p\}.
%     \end{align*}

%     % Reflection
%     Now, as $m$ takes on each residue modulo $p$, so does $-m$, so
%     \begin{align*}
%         p-\nu_A(p) 
%         &= \#\{ m\in\mathbb{Z} : 0\leq m < p \textnormal{ and } -m \not\in |A|_p\} \\
%         &= \#\{ m\in\mathbb{Z} : 0\leq m < p \textnormal{ and } (\forall\ a\in A,\ m+a \not\equiv 0 \pmod p )\}.
%     \end{align*}
    
%     % Grow by p^k
%     For any $m$ satisfying $m + a \not\equiv 0 \pmod p$, $m+jp$ will satisfy that condition for any integer $j$. 
%     So by extending the upper bound of the range condition to $ap$ we have $a$ times as many 
%     values of $m$ satisfying both conditions. Using this with $a=p^{k-1}$ we have
%     \begin{align*}
%         p^k - \nu_A(p)p^{k-1} &=\#\{ m\in\mathbb{Z} : 0\leq m < p^k \textnormal{ and } (\forall\ a\in A,\ m+a \not\equiv 0 \pmod p )\}.
%     \end{align*}

%     Now, $\vk(p)$ contains exactly the integers on the interval $[1,p-1]$, so we may replace the latter condition to obtain
%     $$p^k-\nu_A(p)p^{k-1} = \#\{m\in\mathbb{Z} : 0\leq m < p^k \textnormal{ and } m+A\subset\vk(p)\}.$$
%     Then by Theorem \ref{vk primes} we have
%     \begin{align*}
%         p^k - \nu_A(p)p^{k-1} &= \#\{m\in\mathbb{Z} : 0 \leq m < p^k \textnormal{ and } m+A \subset \vk(p^k) \}\\
%         &= \varphi_A(p^k).
%     \end{align*}
% \end{proof}

\begin{theorem}[Totient Product Formula]~ \\
    For any k-tuple $A$ and natural number $n$, we have:
    $$\varphi_A(n)=n\prod_{p|n}\left(1-\frac{\nu_A(p)}{p}\right).$$
\end{theorem}
\begin{proof}
Let $A$ be a k-tuple and $n$ be a natural number.
Then $n$ has a factorization, say
$$n=\prod_{p_i|n}p_i^{a_i}.$$
Then we have by theorems \ref{phi multiplicative} and \ref{phi prime powers},
\begin{align*}
    \varphi_A(n) 
    &= \varphi_A\left(\prod_{p_i|n}p_i^{a_i}\right) \\
    &= \prod_{p_i|n}\varphi_A(p_i^{a_i}) \\
    &= \prod_{p_i|n} \left(p_i^{a_i} - \nu_A(p_i)p_i^{a_i-1}\right) \\
    &= n\prod_{p_i|n} \left(1 - \frac{\nu_A(p_i)}{p_i}\right).
\end{align*}
\end{proof}

\begin{theorem}
    For any k-tuple $A$, the totient function has the following properties:
    \begin{enumerate}
        \item[\textnormal{(a)}] $\varphi_A(mn) = \varphi_A(m)\varphi_A(n)(d/\varphi_A(d))$, where $d=(m,n)$.
        \item[\textnormal{(b)}] $\varphi_A(mn) = \varphi_A(m)\varphi_A(n)$ if $(m,n)=1$.
        \item[\textnormal{(c)}] $a|b$ implies $\varphi_A(a)|\varphi_A(b)$.
    \end{enumerate}
\end{theorem}
\begin{proof}
    The proof of these follows nearly verbatim the proofs for $\varphi$,
    which is shown in Theorem 2.5 of \cite{Apostol}.
\end{proof}

\begin{definition}
    Let $w$ be an arithmetic function such that:
    \begin{enumerate}
        \item $w$ is completely multiplicative.
        \item $w$ is eventually constant, say $w(p)=k$ for $p>N$.
        \item For all primes $p$, $w(p) \le p$ and $w(p) \le k$.
    \end{enumerate}
    Then we say $w$ is \textbf{unit looking}.
    If there exists a k-tuple such that $\nu_A=w$, we say $A$ generates $w$, and
    say $w$ is good. 
    Otherwise, $w$ is bad.
\end{definition}

Unit looking functions, as far as can be checked by simple methods, appear
to have been generated by some k-tuple.
For $k=2$, all unit looking functions are good.
% Surprisingly, brute force searching indicates for $k>2$ that almost all unit looking 
% functions are bad.
Proving any individual unit looking function is bad is a tedious
affair involving many cases, and so is omitted.
A general algorithm determining this property would be useful in
later computation, but it seems unlikely such an algorithm is feasable
in the general case.
As a concrete example, the unit looking function 
$w$ with $w(2)=1$, $w(3)=w(103)=2$, and $w(p)=3$ for all other primes $p$
is bad, and is the minimal example.
Computation shows that bad unit looking functions become far more common
as the primes at which they differ from their eventual value grow,
and it seems likely that the proportion of good functions goes to zero
in that limit.
We may, however, easily find k-tuples whose unit matches a given unit looking function up 
to an arbitrary bound.

\begin{theorem}
    Let $p_n$ be the $n$th prime number.
    Let $w$ be a unit looking function.
    Then for any natural $N$, there exists a 
    k-tuple $A$ such that $\nu_A(n)=w(n)$ for $n \leq N$.
\end{theorem}
\begin{proof}
    By the complete multiplicativity of $\nu_A$ and $w$,
    we need only show the equality on primes up to $N$.

    Let $k$ be the eventual value of $w$ on primes.
    Let $A_1=\{0\}$.
    If $k=1$, we have $\nu_{A_1}=w$ and are done.
    Now for each $i$ from 2 to $k$, let $Q_i$ be the set of primes up to $N$ for 
    which $\nu_{A_i}(p)<w(p)$.
    By the Chinese Remainder Theorem there exists $a$ which solves $a\equiv i \pmod q$
    for all $q\in Q_i$, and $a\equiv 0$ modulo each other prime up to $N$.
    Let $A_i=A_{i-1}\cup a$.

    Then for each prime $p_n$ up to $N$, exactly the first $w(p_n)$ residues 
    have been filled by elements of $A_k$. 
    That is, $\nu_{A_k}(p_n)=w(p_n)$.
\end{proof}
    
% ==============================================================================================================================
\section{Dirichlet Convolution}
For any arithmetic functions $f$ and $g$ we denote their Dirichlet product by
$$(f\ast g)(n) := \sum_{ab=n}f(a)g(b).$$
We denote the Dirichlet inverse of an arithmetic function $f$ by $f^{-1}$.
We define the arithmetic function $N$ by $N(n)=n$.

\begin{definition}
    For any k-tuple $A$, define $\mu_A$ to be the unique Dirichlet inverse
    of $\nu_A$.
\end{definition}

By definition $\nu_A$ is completely multiplicative, so by Theorem 2.17 
in \cite{Apostol}, we have its Dirichlet inverse to be
$$\mu_A(n)=\nu_A^{-1}(n)=\mu(n)\nu_A(n).$$

We have also from Theorem 2.18 of \cite{Apostol} that for any multiplicative function $f$ and natural number $n$, that
\begin{equation}\label{sum to product}
    \sum_{d|n}\mu(d)f(d) = \prod_{p|n}(1-f(p)),
\end{equation}
which is used to prove the following.

\begin{theorem}
    For any k-tuple $A$, we have
    $$\nu_A*\varphi_A = N.$$
\end{theorem}
\begin{proof}
    Applying equation (\ref{sum to product}) with $f(n)=\nu_A(n)/n$ to the totient product formula, we have
    \begin{align*}
        \varphi_A(n) 
        &= n \prod_{p|n}\left(1-\frac{\nu_A(p)}{p}\right) \\
        &= n \sum_{d|n}\mu(d)\frac{\nu_A(d)}{d} \\
        &= \sum_{d|n}\mu_A(d)\frac{n}{d} \\
        &= (\mu_A\ast N)(n).
    \end{align*}
    Taking the Dirichlet product of both sides by $\nu_A$, we obtain the theorem.
\end{proof}

This result, while simple, is the reason for our later analysis.
Intuitively, $\varphi_A$ should encode information about relative primality
in the k-tuple sense.
The function $N$ is neutral in this regard, so we should expect
$\nu_A=N\ast\varphi_A^{-1}$ to retain some aspect of that information
while playing nicely with Dirichlet series.

\section{Alternatives to $\varphi_A$}
When we defined $\varphi_A$, we counted occurences of complete
coprimality.
We could also ask for the occurences of any amount of coprimality,
a specific amount, or a minimum amount, among other requirements.
However, these examples do not produce multiplicative $\varphi$ analogs.

There is also a subtlety in defining them that doesn't exist in
our $\varphi_A$. 
We defined $\varphi_A$ in terms of translates of $A$, but we could have equivalently
counted the translates of $A_n$.
These definitions are inequivalent for the examples listed above,
for we will have $\#A_n < \#A$ when elements of $A$ coincide modulo $n$.
That is, the issue with those examples is that they rely on the 
number of $a\in k+A$ or $a\in k+A_n$ to satisfy a property, which
may differ. 
We may avoid this issue be only considering functions which
count those $k$ such that each of the elements of $k+A$ possesses
some property, rather considering a property on $k+A$ itself.
But to consider a function to be an analog of $\varphi$, the 
property considered should be ``Are these elements coprime?".
There is only one function of this type which considers
this property: $\varphi_A$.

We may relax our constraints, though, on what property constitutes a totient
to depend on the value of $(k+a,n)$, rather than if it is 1 or not.
This gives not only additional k-tuple totients, but also
additional singular totients.
In general, we may use the property that $(k+a,n)$ is limited to
a specific subset of the primes and obtain a multiplicative totient.
\begin{definition}
    Given a set of primes $B$, let $a \prec B$ mean that $a$ is composed
    only of elements of $B$.
\end{definition}
In particular, if this set is exactly the primes less than some bound $b$,
we are requiring $(k+a,n)$ to be b-smooth:
\begin{definition}
    If the prime factors of an integer $k$ are bounded by $b$, we
    say $k$ is \textbf{$b$-smooth}.
\end{definition}
\begin{definition}
    For a k-tuple $A$ and a set of primes $B$,
    define $\varphi_{A,B}$ by
    \begin{align*}
        \varphi_{A,B}(1) &:= 1, \\
        \varphi_{A,B}(n) &:= \#\{k\in\mathbb{Z}_n : \forall\ a \in A,\
        (k+a,n) \prec B\}.
    \end{align*}
\end{definition}
\begin{theorem}\label{smooth totient mult}
    For any k-tuple $A$ and set of primes $B$,
    $\varphi_{A,B}$ is multiplicative.
\end{theorem}
\begin{proof}
    The proof follows exactly that of $\varphi_A$ and so
    is omitted.
\end{proof}

What, then, are the corresponding unit functions $\nu_{A,B}$?
They are the same as our earlier $\nu_A$, except
they are zero at every number containing a factor
in $B$. 
They are completely multiplicative, can be used to
compute $\varphi_{A,B}$ in the same way, and yield
$\nu_{A,B}\ast \varphi_{A,B} = N$.
The proofs of these facts follow exactly those for $\nu_A$.
We return now to the realm of $\varphi_A$, keeping in mind
that many results may be extended trivially in this way to 
$\varphi_{A,B}$.

The earlier nonmultiplicative examples are not without value,
and not beyond hope.
In essentially all cases, we can express the desired function
as a linear combination of $\varphi_A$. 
Consider first for some $n$ and k-tuple $A$, counting the $k\in\mathbb{Z}_n$
such that exactly $q$ elements of $k+A$ are coprime to $n$.
We call the value $q$ the quota, and refer to this as the unwrapped,
exact quota case.

\begin{definition}[Quota k-tuple totients]
    For any k-tuple $A$ and quota $q\in\mathbb{Z}$ with $0 \le q \le \#A$, define 
    $\varUpsilon^{q=}_A$ by
    \begin{align*}
        \varUpsilon^{q=}_A(1) &:= 1 \\
        \varUpsilon^{q=}_A(n) &:= \#\{k\in\mathbb{Z}_n : q = 
            \#\{a \in A : (k+a,n)=1\} \}.
    \end{align*}
\end{definition}

Replacing $A$ with $A_n$ in this definition yields a wrapped case, and 
requiring at least $q$ values to be coprime to $n$ yields a minimum quota
case. 
We denote the wrapped cases with $\varrho$ rather than $\varUpsilon$,
and minimum quotas use $q\le$ in the superscript.
Clearly, $\varUpsilon_A^{\#A=}=\varphi_A$, though this is also
proven by the following theorem.

\begin{theorem}
    For any k-tuple $A$ and quota $q$,
    $$\varUpsilon^{q=}_A = \sum_{k=q}^{\#A}(-1)^{k-q}
        { q \choose k }
        \sum_{\substack{B\subset A \\ \#B = k}}\varphi_B$$
\end{theorem}
\begin{proof}
    TODO
\end{proof}

\begin{theorem}
    For any k-tuple $A$, and quota $q$,
    $$\varUpsilon^{q\le}_A = \sum_{k=q}^{\#A}(-1)^{k-q}
        {k-1 \choose k-q}
        \sum_{\substack{B\subset A \\ \#B = k}}\varphi_B.$$
\end{theorem}
\begin{proof}
TODO
\end{proof}

The formulae for the corresponding $\varrho$ functions are identical:
simply replace $A$ with $A_n$ throughout.
This creates a dependance on $n$ in the final sum which appears to
explode the computational complexity, but for sufficiently
large $n$, $A = A_n$, so the order of computational complexity is the same.
This also shows that the wrapped and unwrapped totients are
eventually equivalent, as should be expected.
We will make no further use of the wrapped totients, and the proofs are
mostly the same as the unwrapped cases, so they are omitted.

Using the discrete modular convolution, we can obtain
the essential data required to compute wrapped quota totients
for all quotas at once.

\begin{definition}
    Given $f,g : \mathbb{Z}_n \rightarrow \mathbb{C}$,
    define the discrete modular convolution $\odot_n$ by
    $$(f\odot_n g)(k) = \sum_{\substack{a+b\equiv k \\ \bmod n}}f(a)g(b),$$
    where $a,b\in\mathbb{Z}_n$. 
\end{definition}

The symbol $\odot$ is non-standard. 
Normally $\ast$ would be used, but it is already in use for Dirichlet convolution.

\begin{theorem}
    Let $A\in\mathbb{K}$,
    let $n\in\mathbb{N}$, and let $\chi_0$ be the principle 
    Dirichlet character mod $n$.
    Let $1_{A_n}$ be the indicator function of $A_n$.
    Then
    $$(1_{A_n}\odot_n\chi_0)(k) = \#\{a \in A_n : (k+a,n)=1\}.$$
\end{theorem}
\begin{proof}
    TODO
\end{proof}

Computing $1_{A_n}\odot_n\chi_0$ for all $k$ then gives a convenient data
format for querying quota totient information at $n$.
This can be made to work for unwrapped quota totients by
considering a period of size $n$ of the non-modular discrete convolution
of $\chi_0$ and $1_A$.

% Another convolution centric idea, we have $\varphi_A(n)=\#(\varkappa_n +_n A)$

% =======================================================================================================================================
\section{Admissibility}

\begin{definition}
    We say a k-tuple $A$ is \textbf{admissible} if, for each prime $p$, we have
    $$\nu_A(p)<p.$$
\end{definition}

Then a k-tuple is admissible if it does not fill all of the residue classes modulo any prime.
This condition gives a simple bound on $\nu_A$.

\begin{proposition}
    For any admissible k-tuple $A$, we have
    $$\nu_A(n) \leq \varphi(n).$$
\end{proposition}
\begin{proof}
    Let $A\in\mathbb{K}$ be admissible, and let $n$ be a natural number with factorization
    $$n=\prod_{p_i|n}p_i^{a_i}.$$
    Then 
    \begin{align*}
        \nu_A(n) 
        &= \nu_A\left(\prod_{p_i|n}p_i^{a_i}\right) \\
        &= \prod_{p_i|n}\nu_A(p_i)^{a_i} \\
        &\leq \prod_{p_i|n}(p_i-1)^{a_i} \\
        &\leq \prod_{p_i|n}(p_i-1)\cdot p_i^{a_i-i} \\
        &= \varphi(n).
    \end{align*}
\end{proof}

We are now equipped to state the prime k-tuple conjecture.

\begin{conjecture}[Prime K-Tuple Conjecture]
For any finite, admissible k-tuple $A$, there exists an infinite k-tuple $B$ such that $A+B\subset \mathbb{P}$.
\end{conjecture}

Admissibility is required, for if a k-tuple $A$ was inadmissible by prime $p$, for $b\in B$ large enough that each element
of $b+A$ is greater than $p$, we must have some element of $A$ filling the 0 residue.
That element is then divisible by $p$, $b+A\not\subset\mathbb{P}$, and it cannot be for any $b$ beyond that bound.
Thus, there can be only finitely many integers at which $b+A\subset\mathbb{P}$ and no valid $B$ exists.

The finiteness of $A$ is also required. 
For example, let $A$ be the set of all primorials.
This fills at most $i$ residues of each prime $p_i$ and so is admissible.
We cannot have negatives in $B$, for $2\in A$ is the minimum prime.
We cannot have 0 or 1 in $B$, because $A+B$ would contain 6 and 4 respectively.
If $B$ contains $b>1$, $b$ is divisible by some prime $p$, and the primorials will eventually be divisible by $p$.
Then for some element $a\in A$ divisible by $p$, $p|a+b$, so $A+B\not\subset\mathbb{P}$.
Thus, there are no integers $b$ such that $b+A\subset\mathbb{P}$.
Not only can $B$ not be infinite in this case, it cannot contain any integer.

If the conjecture holds, there exist infinite k-tuples $A$ with arbitrary large finite $B$ such that
$A+B\subset\mathbb{P}$. 
Indeed, given a finite admissible k-tuple $B$, we apply the conjecture and find such an $A$.
The question of adding two infinite k-tuples and obtaining a prime k-tuple is handled in greater generality in 
\cite{Granville},
the results of which we state in less generality as the following theorem. 
It is often said that the Prime K-Tuple Conjecture feels intuitively true, so it is somewhat alarming 
that it holds as a consequence the following result, which seems intuitively impossible.
\begin{theorem}
    Assume the Prime K-Tuple Conjecture.
    Then there exist infinite k-tuples $A$ and $B$ such that $A+B\subset\mathbb{P}$.
\end{theorem}

% ========================================================================================================================
% ========================================================================================================================
% ========================================================================================================================

% The function $\nu_A$ is convenient to work with in the setting of arithmetic functions in
% that it takes on only integer values and gives the relationship $\nu_A \ast \phi_A = N$.
% Transporting it to the domain of Dirichlet series, though, poses a challenge.
% The function $\mathcal{D}(\nu_A)(s)=\sum_{n\ge 1}\nu_A(n)n^{-s}$ permits an Euler product, 
% and that product can be written as a power of $\zeta$, a finite product, and 
% the series $\Di((\#A)^{\Omega})$, where $\Omega(n)$ is the number of primes 
% (counting multiplicity) which compose $n$.
% While the final portion of this form has reasonable convergence properties,
% it will be the cause of much of the complexity in analyzing $\nu_A$.
% That portion, though, does not seem to have much to do with the additive patterns
% we seek to analyze; rather it describes that $\nu_A$ should grow
% approximately with the number of primes in its composition.
% The following adjustment, then, gives an arithmetic function
% which contains the same essential information, but which has far simpler
% analytic properties.

% \begin{definition}
%     Let $A$ be a k-tuple.
%     For each $n$, define $w_A(n)$ by
%     $$w_A(n) = \frac{\nu_A(n)}{(\#A)^{\Omega(n)}}.$$
% \end{definition}

% This is equivalent to defining $w_A(p)=\nu_A(p)/\#A$ on primes and
% extending it by complete multiplicativity as with $\nu_A$,
% but we may also show this directly: because $\Omega$ is completely additive, 
% $(\#A)^{\Omega}$ is completely multiplicative, so $w_A$ is as well.
% Note also that for primes $p$ with $\max(A) \le p$, $\nu_A(p)=\#A$, so
% $w_A$ is eventually 1 on primes.
% We will work with both $\Di(\nu_A)$ and $\Di(w_A)$.

% \section{The Dirichlet Series of $w_A$}
 
% We have that $w_A$ is completely multiplicative, so $\Di(w_A)$ permits the Euler product
% $$\Di(w_A)(s) = \prod_{p}(1-w_A(p)p^{-s})^{-1},$$
% but we know for sufficiently large primes, $w_A(p)=1$, so
% letting $S_A$ be the set of primes on which $w_A(p)\not= 1$, we have
% $$\Di(w_A)(s) = \zeta(s)\prod_{p\in S_A}\frac{1-p^{-s}}{1-w_A(p)p^{-s}}.$$
% Let $M_A(s)$ denote the finite product above and call this the modulator for $A$.
% This form immediately yields the analytic continuation of $\Di(w_A)$ via
% the continuation of $\zeta$ to the entire plane, except at the pole contributed 
% by $\zeta$ at $s=1$, and one infinite (but well behaved) set of poles for each $p\in S_A$.
% These poles are at the solutions of $0=1-w_A(p)p^{-s}$, which are found, for any 
% $k\in\mathbb{Z}$, at $s_{p,k} = \frac{\log(w_A(p)) + 2\pi k i}{\log(p)}$.

% \begin{proposition}
%     For any admissible k-tuple, we have 
%     $$\frac{1}{x}\sum_{n\le x} w_A(n) \sim M_A(1).$$
% \end{proposition}
% \begin{proof} 
%     Every modulator pole has real part less than one.
%     Thus, we can apply the Wiener-Ikehara theorem (1.1 in \cite{Korevaar}) 
%     which gives
%     $$\frac{1}{n}\sum_{n\le x} w_A(n) \sim \text{Res}(\Di(w_A), 1).$$
%     We have $\text{Res}(\zeta,1)=1$, so
%     $$\text{Res}(\Di(w_A), 1) = M_A(1),$$
%     as was needed.
% \end{proof}

% To find the residues of any pole of the modulator, we first 
% show that no two poles among those listed coincide. 
% Clearly none of those given by the modulator coincide with $s=1$, so we need only 
% check that no poles among the modulator coincide.
% Suprisingly, though, this will walk us directly into the four exponentials conjecture.
% \begin{conjecture}[Four exponentials conjecture]
%     For $1\le i,j \le 2$, let $\lambda_{ij}\in\mathbb{C}$ such that
%     each $exp(\lambda_{ij})$ is algebraic. Suppose that $\lambda_{11}$
%     and $\lambda_{12}$ are linearly independent, and that $\lambda_{21}$
%     and $\lambda_{22}$ are linearly independent, both in the rationals.
%     Then
%     $$\lambda_{11}\lambda_{22} \not= \lambda_{12}\lambda_{21}.$$
% \end{conjecture}
% \begin{lemma}
%     Assume the four exponentials conjecture. 
%     Then we have
%     $$s_{p_1,k_1} = s_{p_2,k_2} \implies (p_1,k_1) = (p_2,k_2).$$
% \end{lemma}
% \begin{proof}
%     % Let $a_i = w_A(p_i)$ and note that $a_i\in\mathbb{Q}$ and $0 < a_i < 1$, 
%     % so $a_i\not\in\mathbb{Z}$.
%     Assume $s_{p_1,k_1} = s_{p_2,k_2}$, and assume for contradiction $p_1\not=p_2$
%     Then their real parts coincide and we have $\log{a_1}/\log{p_1} =\log{a_2}/\log{p_2}$,
%     which gives 
%     \begin{equation}\label{four exp 1}
%     \log{a_1}\log{p_2} = \log{a_2}\log{p_1}.
%     \end{equation}
%     Each $a_i$ and $p_i$ is rational and thus algebraic, so the exponential
%     of each logarithm is algebraic.
    
%     To show the linear independence of $\log{p_1}$ and $\log{p_2}$, assume
%     there are $r,s\in\mathbb{Q}$ such that $r\log{p_1} + s\log{p_2} = 0$.
%     Then $p_1^rp_2^s = 1$, which is only possible if $r=s=0$.

%     If $\log{a_1}$ and $\log{a_2}$ are dependent, then there are $r,s\in\mathbb{Q}$,
%     not both 0, such that $r\log{a_1} + s\log{a_2} = 0$.
%     Without loss of generality, assume $r\not=0$.
%     Then $a_1 = a_2^{-s/r}$.
%     Applying this to equation (\ref{four exp 1}), we have
%     \begin{align*}
%     -\frac{s}{r}\log{p_2} &= \log{p_1} \\
%     \implies p_2^{-s/r} &= p_1,
%     \end{align*}
%     which is a contradiction.
%     Thus, $\log{a_1}$ and $\log{a_2}$ are linearly independent.

%     Finally, we may apply the four exponentials conjecture, which 
%     contradicts (\ref{four exp 1}).
%     Thus, $p_1=p_2$. 
%     Showing $k_1=k_2$ is trivial.
% \end{proof}

% The values of $a_i$ in the above proof are not entirely arbitrary, coming
% from $w_A$. 
% As such, there may be some way of showing this result holds
% without the four exponentials conjecture.
% We continue to assume the conjecture as we now find the residues of the poles
% associated with the modulator.

% \begin{theorem}
%     $$\text{Res}(\Di(w_A), s_{p,k}) = 
%         \frac{1-p^{-s_{p,k}}}{\log{p}} \zeta(s_{p,k}) 
%         \prod_{\substack{q\in S\\q\not=p}}\frac{1-q^{-s_{p_k}}}{1-w_A(q)q^{-s_{p,k}}}.
%         $$
% \end{theorem} 
% \begin{proof}
%     By the lemma, each pole of $\Di(w_A)$ is simple.
%     That is, for any $s_{p,k}$ we have $\Di(w_A)(s) = 
%     H_{p,k}(s)\cdot g(s)$, where 
%     $$g(s) = \frac{1}{1-w_A(p)p^{-s}},$$
%     and
%     $$H_{p,k}(s) = \zeta(s)(1-p^{-s})
%         \prod_{\substack{q\in S\\q\not=p}}\frac{1-q^{-s_{p_k}}}{1-w_A(q)q^{-s_{p,k}}},
%         $$ 
%     where $H_{p,k}$ is holomorphic at $s_{p,k}$.
%     Now, note that by the definition of $s_{p,k}$, $1-w_A(p)p^{-s_{p,k}}=0$.
%     Using this, we have
%     \begin{align*}
%     \text{Res}((1-w_A(p)p^{-s})^{-1},s_{p,k}) 
%         &= \frac{1}{\frac{d}{ds}(1-w_A(p)p^{-s})|_{s_{p,k}}} \\ 
%         &= \frac{1}{w_A(p)\log(p)p^{-s_{p,k}}} \\ 
%         &= \frac{1}{\log(p)},
%     \end{align*}
%     which gives
%     \begin{align*}
%         \text{Res}(\Di(w_A),s_{p,k}) = \frac{H_{p,k}(s_{p,k})}{\log(p)}
%     \end{align*}
%     as was needed.
% \end{proof}


% ================================================================================
\section{The First Hardy-Littlewood Conjecture}

\begin{definition}
    For any k-tuple $A$, define 
    $$\pi_A(x) = \#\{n \le x : n + A \subset \mathbb{P} \}.$$
\end{definition}

\begin{definition}
    For $n>0$, define 
    $$\Li_n(x) = \int_{2}^{\infty}\frac{\mathrm{dt}}{\log^n(t)}.$$
    We call this the $n$th (offset) logarithmic integral.
\end{definition}
% The ``offset" is due to the lower bound being 2, rather than 0.
% The prime number theorem is usually given in the form
% $\pi(x) \sim \frac{x}{\log(x)}$,
% but we have also that $\Li_1(x) \sim \frac{x}{\log(x)}$, and it is 
% well known that $\Li_1$ gives a better approximation of $\pi$.
% It is no surprise then that in 1923, Hardy and Littlewood made
% the following conjecture.
\begin{conjecture}[First Hardy-Littlewood conjecture]\label{HLC}
Let $A$ be an admissible k-tuple. Then
$$\pi_A(x) \sim C_A\Li_{\#A}(x),$$
where
$$C_A = \prod_{p}\frac{1-\nu_A(p)p^{-1}}{(1-p^{-1})^{\#A}}.$$
\end{conjecture}

This conjecture can be understood intuitively.
Consider a translation of a k-tuple $A$ by a large value $j$.
If we assume that prime numbers are distributed randomly with 
a distribution of $1/\log(n)$, and that the probabilities of each element
of $A+j$ being are uncorrelated, we would guess the probability of all the 
elements of $A+j$ being prime to be
$$\prod_{n\in j+A}1/\log(n).$$
For our large $j$, this is essentially $1/\log(n)^{\#A}$.
Then we would naively expect $\pi_A(x) \sim \Li_{\#A}(x)$.

The presence of $C_A$, refered to as the `singular series' for $A$,
is a correction factor for our assumption of non-correlation.
To see this, consider $A=\{0,2\}$, the twin prime case.
Translating this by our large $j$, if $j+0$ is prime,
we know that $j+0$ is odd. 
Then $j+2$ is also odd, which should double the chance that $j+2$ is prime!
We have also that $j+0\equiv 1,2 \pmod 3$, so $j+2\equiv 0,1 \pmod 3$.
Our assumption that primes behave randomly essentially assumes that
there was a $1/3$ chance that $j+2$ was a multiple of 3 
regardless of the primality of $j+0$; we now see that the structure of $A$
dictates the actual chance should be $1/2$.
Correcting probabilities with respect to the structure of $A$ for each 
prime in this way gives exactly $C_A$.
This sort of probabilistic game allows the guessing of many
correct formulae for prime numbers.

This conjecture, of course, implies the prime k-tuple conjecture.
It also implies that the exact structure of a k-tuple
does not determine its asymptotic behavior, but
rather that behavior is entirely decided by $\nu_A$.
For example, $\{0,2\}$ and $\{0,4\}$ produce the same $\nu_A$.
This yields an equivalence relation.

\begin{definition}
    Given k-tuples $A$ and $B$, we say they are \textbf{residually equivalent} if
    $\nu_A=\nu_B$.
\end{definition}

The following asymptotic expansion is well known:
\begin{equation}\label{Li1 asymp}
\Li_1(x) \sim \frac{x}{\log(x)}\sum_{k\ge 0} \frac{k!}{\log^{k}(x)}.
\end{equation}
The simpler estimate $\frac{x}{\log{x}}$ given in the prime number theorem is then
the truncation of this expansion to its first term.
Then we should expect to analagously use the truncation of 
the asymptotic expansion of $\Li_n$ in attempting to gain traction in the 
first Hardy-Littlewood conjecture.

\begin{proposition}
    TODO: this result needs additional care due to the meaning of asymptotic
    expansion, though it is correct.
    We could also just state it as true, I don't think proving this
    is all that important.
    We have the asymptotic expansion of $\Li_n$ to be
    $$\Li_n(x) \sim \frac{x}{(n-1)!\log{x}}\sum_{k\ge n-1}\frac{k!}{\log^k(x)}.$$
\end{proposition}
\begin{proof}
    First, apply integration by parts to $\Li_n$ with $dv = dx/(x\log^n(x))$ and
    $u=x$ to obtain
    \begin{equation}\label{Lin asymp ind}
        \Li_n(x) = \frac{-x}{(n-1)\log^{n-1}{x}} + \frac{1}{n-1}\Li_{n-1}(x)
    \end{equation}
    for $n>1$.
    We have the result for $n=1$: it is exactly equation (\ref{Li1 asymp}).
    Assume then for induction the result holds for some $n$.
    Then applying equation (\ref{Lin asymp ind}) we have
    \begin{align*}
        \Li_{n+1}(x) &= \frac{-x}{n\log^{n}{x}} + \frac{1}{n}\Li_{n}(x) \\
        &\sim \frac{-x}{n\log^{n}{x}} 
            + \frac{x}{n!\log{x}}\sum_{k\ge n-1}\frac{k!}{\log^kx} \\
        &= \frac{x}{n!\log{x}}\left(\frac{-(n-1)!}{\log^{n-1}{x}} +
            \sum_{k\ge n-1}\frac{k!}{\log^k{x}}\right) \\
        &= \frac{x}{n!\log{x}}\sum_{k\ge n}\frac{k!}{\log^k{x}}.
    \end{align*}
\end{proof}

All this is to say: if the first Hardy-Littlewood conjecture is correct,
we should attempt to show, following in the steps of the prime number theorem,
that 
$$\pi_A(x) \sim \frac{C_Ax}{\log^{\#A}{x}}.$$

\section{The Dirichlet Series of $\nu_A$}

We denote the Derichlet series of an arithmetic function $f$ by
$$\Di(f)(s) := \sum_{n\ge 1}\frac{f(n)}{n^{s}}.$$

% For each admissible constellation $A$, Conjecture \ref{HLC} gives an associated 
% constant $C_A$.
% This section aims to explore these constants in our context.
% Just as we describe $\Di(w_A)$ in terms of $\zeta$ and a finite modulator, 
We may factor $\Di(\nu_A)$ as
\begin{equation}\label{DUA zeta form}
\Di(\nu_A)(s) = \zeta^{\#A}(s)  
    \prod_{p}\frac{(1-p^{-s})^{\#A}}{1-\nu_A(p)p^{-s}}.
\end{equation}

We refer to the final product as the modulator, denoted $N_A(s)$.
The following result essentially states that the singular series
$C_A$ converges, which is well known, but is given for completeness.

\begin{proposition}
    For any admissible k-tuple $A$, there exists $\epsilon > 0$
    such that $N_A(s)$ converges for $\Re(s) > 1-\epsilon$.
\end{proposition}
\begin{proof}
    Let $A$ be an admissible k-tuple.
    Then $\nu_A(p) < p$, so the real parts of the solutions of $1-\nu_A(p)p^{-s}=0$ are all
    left of 1, and go to 0 as $p\rightarrow\infty$.
    Then there exists $\epsilon>0$ such that all solutions of $1-\nu_A(p)p^{-s}=0$
    lie to the left of $1-\epsilon$.
    Fix $s=\sigma + it$ with $\sigma>1/2$.
    Let 
    $$N_{A,r}(s) = \prod_{p>r}\frac{(1-p^{-s})^{\#A}}{1-\nu_A(p)p^{-s}}.$$
    We will show the convergence of $N_{A,r}$ at $s$ for a fixed $r$, from which 
    the result follows, as the finitely many removed factors all converge
    for $\Re(s)>1-\epsilon$.
    We fix $r$ to be large enough that both $\#Ar^{-\sigma} < 1/2$,
    and $\nu_A(p)=\#A$ for $p>r$.
    We then have
    \begin{align*}
        \log{N_{A, r}}(s) 
        &= \log{\prod_{p>r}\frac{(1-p^{-s})^{\#A}}{1-\#Ap^{-s}}} \\
        &= \sum_{p>r}\left(\#A\log(1-p^{-s}) - \log(1-\#Ap^{-s})\right).
    \end{align*}
    Expanding the logs into power series gives
    \begin{align*}
        \log{N_{A,r}}(s) 
        &= \sum_{p>r}\left(
            -\#A\sum_{k\ge 1}\frac{p^{-sk}}{k} + 
            \sum_{k\ge 1}\frac{(\#Ap^{-s})^{k}}{k}\right) \\
        &= \sum_{p>r}\sum_{k\ge 2}\left(
            \frac{(\#A)^{k}-\#A}{k}p^{-sk}
            \right).
    \end{align*}
    Note the cancellation of the $k=1$ term.
    Now,
    \begin{align*}
        |\log{N_{A,r}}(s)| 
        &\le \sum_{p>r}\sum_{k\ge 2}
            \left(\#Ap^{-\sigma}\right)^k.
    \end{align*}
    For each $p>r$, by our bound on $r$, the $k$ sum is a convergent geometric series.
    We have also for $p>r$ that $1-\#Ap^{-\sigma} > 1/2$, so
    \begin{align*}
        |\log{N_{A,r}}(s)| 
        &= \sum_{p>r}\frac{(\#Ap^{-\sigma})^2}{1-\#Ap^{-\sigma}} \\
        &\le 2(\#A)^2\sum_{p>r}p^{-2\sigma} \\
        &\le 2(\#A)^2\sum_{n>r}n^{-2\sigma}.
    \end{align*}
    We have $\sigma > 1/2$, so this is a convergent p-series.
    Thus, $\log{N_{A,r}}$ converges absolutely, and exponentiating gives
    the result.
\end{proof}

We do not need the full domain of convergence of the prior result; our
goal was to justify evaluating $N_A(1)$ and obtaining the
following.

\begin{corollary}
    For any admissible k-tuple $A$, we have
    $$\frac{1}{C_A} = N_A(1),$$
    and
    $$\Res(\zeta^{\#A},1) = C_A\Res(\Di(\nu_A),1).$$
\end{corollary}
\begin{proof}
    The first equality is immediate from the definitions of $C_A$
    and $N_A$.
    Taking residues in (\ref{DUA zeta form}) then yields the second equality.
\end{proof}

The Laurent expansion of $\zeta$ is known, so the residues above
are computable.
We may factor $N_A$ further into $\Di((\#A)^\Omega)$ and
a finite product, but for k-tuples with a prime size, this creates
a pole and zero at $s=1$ in those factors respectively.
% Still, considering $\Di((\#A)^\Omega)$ contains no information
% on the structure of $A$ outside of its size, we should anticipate
% the finite product in that factorization to possess some structural information
% akin to $M_A$.
We now compute $N_A(1)$ for some 
admissible, non-residually-equivalent k-tuples,
where the $p$ column gives those primes over which $\nu_A(p)\not=\#A$.
% Computations for $M_A$ were made entirely in the rationals before
% rounding and so should be accurate.
Computations were made using the product over the first 
10000 primes before rounding, and are made only for the sake of discussion:
they have not been verified for accuracy.
We also include the corresponding $C_A$s.

\begin{center}
% \begin{tabular}{ c | c | c | c | c }
%     $A$ & $p$ & $M_A(1)$ & $N_A(1)$ & $C_A$ \\ 
%     \hline
%     $\{0,2\}$         & 2          & 0.66666667 & 0.75738949 &  1.32032464 \\
%     $\{0,10\}$        & 2,5        & 0.59259259 & 0.56804211 &  1.76043285 \\
%     $\{0,6\}$         & 2,3        & 0.53333333 & 0.37869474 &  2.64064928 \\
%     $\{0,2,6,8\}$     & 2,3        & 0.45714286 & 0.24089421 &  4.15119988 \\
%     $\{0,8,12,30\}$   & 2,3,5,11   & 0.41975814 & 0.10539121 &  9.48845686 \\
%     $\{0,6,12,18\}$   & 2,3        & 0.41558441 & 0.12044710 &  8.30239975 \\
%     $\{0,2,6,8,30\}$  & 2,3,5,7,11 & 0.38614839 & 0.05639911 & 17.73077650 \\
%     $\{0,18,30,60\}$  & 2,3,5,7    & 0.35463203 & 0.03011178 & 33.20959900 \\
%     $\{0,60,120,180\}$& 2,3,5      & 0.34996582 & 0.03011178 & 33.20959900 \\
% \end{tabular}
\begin{tabular}{ c | c | c | c }
    $A$ & $p$ & $N_A(1)$ & $C_A$ \\ 
    \hline
    $\{0,2\}$         & 2          & 0.75738949... &  1.32032464... \\
    $\{0,10\}$        & 2,5        & 0.56804211... &  1.76043285... \\
    $\{0,6\}$         & 2,3        & 0.37869474... &  2.64064928... \\
    $\{0,2,6,8\}$     & 2,3        & 0.24089421... &  4.15119988... \\
    $\{0,8,12,30\}$   & 2,3,5,11   & 0.10539121... &  9.48845686... \\
    $\{0,6,12,18\}$   & 2,3        & 0.12044710... &  8.30239975... \\
    $\{0,2,6,8,30\}$  & 2,3,5,7,11 & 0.05639911... & 17.73077650... \\
    $\{0,18,30,60\}$  & 2,3,5,7    & 0.03011178... & 33.20959900... \\
    $\{0,60,120,180\}$& 2,3,5      & 0.03011178... & 33.20959900... \\
\end{tabular}
\end{center}

% The table is sorted by $M_A(1)$, and at first glance it appears
% that $N_A(1)$ follows in step.
% The entries $\{0,8,12,30\}$ and $\{0,6,12,18\}$ have been added
% to dispel this. 
% TODO: I don't think this next thing is really worth saying if we don't have 
% anything else to say about it.
% Indeed, we have for k-tuples $A$ and $B$ with $\#A = \#B = c$,
% \begin{align*}
%     M_A(1) &< M_B(1) \\
%     \iff 1 &< \prod_{p}\frac{1-w_A(p)p^{-1}}{1-w_B(p)p^{-1}} \\
%     \iff 0 &< \sum_{p}\left[ \log(1-w_A(p)p^{-1}) - \log(1-w_B(p)p^{-1}) \right] \\
%     \iff 0 &< \sum_{p}\sum_{k\ge1} \frac{w_B(p)^k-w_A(p)^k}{kp^k} \\
%     \iff 0 &< \sum_{k\ge1} \frac{1}{k}\sum_{p} \frac{w_B(p)^k-w_A(p)^k}{p^k}. \\
% \end{align*}
% Following the same proceedure for $N_A(1)$ gives
% \begin{align*}
%     N_A(1) &< N_B(1) \\
%     \iff 0 &< \sum_{k\ge1}\frac{c^k}{k}\sum_{p}\frac{w_B(p)^k-w_A(p)^k}{p^k}, \\
% \end{align*}
% so the ordering of $N$s is essentially that of the $M$s, but with a higher weight
% given to higher powers of each prime.

We see that the final two entries appear to be exactly
equal. It is not a rounding error: they are identical.
This, along with their corresponding k-tuples having the same size,
means that the Hardy-Littlewood Conjecture predicts these k-tuples
to have identical limiting frequency, despite their residual
inequivalence.

When k-tuples $A$ and $B$ have the same size, the ratio $N_A(s)/N_B(s)$ 
is a finite product, so $N_A(1)/N_B(1)$ is rational.
It is convenient to handle these constants in a rational form, so
we would like to define, for each $k$, a constant representing
the infinite portion of the product.
The naive candidate is $\prod_{p}\frac{(1-p^{-1})^{\#A}}{1-p^{-s}}$,
but this diverges.
We could use $N_A(1)$ for any particular admissible $A$, and opt to choose
the maximum among such constants.
\begin{definition}
    For each $k\in\mathbb{N}$, define 
    $$\mathfrak{U}_k = \max\{N_A(1) : \#A = k\}.$$
\end{definition}

\begin{proposition}
    For any $k\in\mathbb{N}$, we have
    $$\mathfrak{U}_k = N_A(1)$$
    for $A$ containing ???
\end{proposition}
\begin{proof}
    I think it is OEIS: A135311.
    Verify via simulation.
    I don't know that this is really important? 
    We don't need to choose one, and also where would we even
    go with this?
\end{proof}

Some of these minimal constants are given below.
\begin{center}
\begin{tabular}{ c | c }
    $k$ & $\mathfrak{U}_k$ \\ 
    \hline
    1 & 1.00000000... \\
    2 & 0.75738948... \\
    3 & 0.34986379... \\
    4 & 0.24089420... \\
    5 & 0.09869844... \\
    6 & 0.05780744... \\
    7 & 0.01852784... \\
    8 & 0.00560960... \\
    9 & 0.00158709...
\end{tabular}
\end{center}


\section{Proof of the First Hardy-Littlewood Conjecture}

Sorry, just kidding.
What, though, are the difficulties in proving such a result?
We've seen that many results may be extended to k-tuples
without much difficulty, so why can we not apply the logic
used in the Prime Number Theorem to its k-tuple analog, 
the First Hardy-Littlewood Conjecture?

The Prime Number Theorem is usually proven in the equivalent form
$\Psi \sim x$, where $\Psi(x)=\sum_{n\le x}\Lambda(n)$
and $\Lambda(n)$ is defined on powers of primes by 
$\Lambda(p^k)=\log(p)$, and is zero otherwise.
The k-tuple analog of $\Psi$, $\Psi_A$, is known in that
the we can obtain the desired result from $\Psi_A \sim C_Ax$.
That function is defined by
$$\Psi_A(x) = \sum_{n\le x}\prod_{a\in A+n}\Lambda(a).$$

We have tools for analyzing summatory functions using
Dirichlet series. 
In particular, $\Psi$ and $\Di(\Lambda)$ are closely related,
and one can show that $\Di(\Lambda)$ is the negative logarithmic
derivitive of $\zeta$.
From this point the Prime Number Theorem is reached through
a careful analysis.
The difficulty in extending this result to k-tuples seems
to lie
in the incompatibility between Dirichlet series, pointwise
products, and shifted sequences.
That is, we would need to work with Dirichlet series of
the form
$$\sum_{n\ge 1} \frac{f(n)g(n+k)}{n^s},$$
and 

It seems that the difficulty in extending this idea lies
in TODO

Testing on general Dirichlet series, we can get the summatory of $\Psi_A$
in terms of a GDS with $\lambda_n = \prod_{a\in A}\log(a+n)$ 
using general Perron's formula.
Let $A^s := e^{s\prod_{a\in A}\log(a)}$.
Then the GDS takes the form
$$\Di_A(\Lambda_A)(s) = \sum_{n}\frac{\Lambda(A+n)}{(A+n)^s},$$
and Perron's gives
\begin{align*}
    \Psi_A(x) 
    &= \sum_{n\le x}\Lambda(A+n) \\
    &= \sum_{\lambda_n \le \log(A+x)}\Lambda(A+n) \\
    &= \frac{1}{2\pi i}\int_{c-i\infty}^{c+i\infty}\Di_A(\Lambda_A)(s)
        \frac{(A+x)^s}{s}ds
\end{align*}
From there you can follow the PNT path up to
$$\Di_A(\Lambda_A)(s) = -\frac{d}{ds}\sum_{\substack{n \\ n + a_i = p_i^{k_i}}}
    \frac{(e^{-s\log(p_1)\cdot\cdot\cdot\log(p_m)})^
        {k_1\cdot\cdot\cdot k_m}}{k_1\cdot\cdot\cdot k_m}.$$
In PNT we just have n to be a power of a prime so we can sum
over all primes of sum over all powers, and the power sum is
a log expansion.
Numeric test shows we don't have 
$$\frac{\zeta_A'}{\zeta_A} = -\Di_A(\Lambda_A)$$
for these tuple Dirichlet series (tds) as we do with singular,
but there may be a more complex link between the TDS of $\Lambda_A$
and the TDS zeta.

Another alternative is multivariable dirichlet series of the form
$$\Di_A(f)(s_1,...,s_k) = \sum_{n}\frac{f(n)}
    {(n+a_1)^{s_1}\cdot\cdot\cdot(n+a_m)^{s_m}}.$$
We'll need to look into modifying Perron's formula to work with this
format, but it seems probable it can be done.
This series lets us do
\begin{align*}
    \Di_A(\Lambda_A)(s_1,...,s_k) 
    &= \sum_{n}\frac{\prod_{a\in A}\Lambda(a+n)}{(n+a_1)^{s_1}\cdot\cdot\cdot(n+a_m)^{s_m}} \\
    &= -\frac{d}{ds_1}\sum_{\substack{n\\n+a_1 = p^k}}
        \frac{p^{-ks_1}\prod_{\substack{a\in A \\ a \not= a_1}}\Lambda(a+n)}
        {k(n+a_2)^{s_1}\cdot\cdot\cdot(n+a_m)^{s_m}}.
\end{align*}
In other words, we can manipulate single logs at a time through derivatives,
rather than having to pop out all logs at once as in the avove case.
This faces the same difficulties, though, where $(n+a_i)^{s_i}$ is
then nonconstant wrt $k$, so even if we split the $n$ sum into $p$ and
$k$ sums, we can't collapse as a log expansion.

There's another way to express the multivariable type: for twin primes it's
$$\mathcal{D}_{\{0,2\}}(\Lambda)(s_1,...,s_m) = \sum_{n_1}\sum_{n_2}\frac{\Lambda(n_1)\Lambda(n_2)}{n_1^{s_1}n_2^{s_2}}\cdot 1_{n_2=n_1+2}.$$
That is, it's the previous one but you expand the sum to every
option and then filter down.
You can get this to
\begin{align*}
& \sum_{n}\sum_{a+b=n}
    \frac{\Lambda(a)\cdot 1_{n=2a+1}}{a^{s_1}}
    \frac{\Lambda(b)}{b^{s_2}} \\
= &\sum_{n}
    \left[\frac{\Lambda\cdot 1_{n=2N+1}}{N^{s_1}}
    *_+
    \frac{\Lambda}{N^{s_2}}\right](n).
\end{align*}

The ability to express this Dirichlet type series as a linear sum
over Cauchy products is suggestive that it might be used on a larger
category of mixed (add + mult) problems such as Bateman-Horn (though
I imagine this corresponds to an even broader scope?).
I think the first generalization, TDS, is the best fit for HLC1,
but the Dirichlet series for $\nu_A$ still needs to get involved somewhere
along the way.
It also faces the ever looming problem of shift coherence, rather than
shift invariance.
The jump between coherent and invariant objects seems critical.

% =======================================================================================================================================
\pagebreak
\section{TODO}

% We may express the limiting behavior of the summatory function of $\nu_A$ via
% $N_A(1)$.
% \begin{theorem}
%     For any admissible k-tuple $A$, we have
%     $$\frac{1}{x\log^{\#A-1}{x}}\sum_{n\le x}\nu_A(n) \sim \frac{N_A(1)}{(\#A-1)!}.$$
% \end{theorem}
% \begin{proof}
%     This is exactly the result of Theorem 5.2 in \cite{Tenenbaum}. 
% \end{proof}

% To that end, we follow the usual route and define:
% \begin{definition}
%     $$\upvartheta_A(x) = \sum_{\substack{n\le x \\ n + A \subset 
%         \mathbb{P}}}\log^{\#A}{n}.$$
% \end{definition}

% \begin{proposition}
%     For $x\ge2$ we have
%     $$\upvartheta_A(x) = \pi_A(x)\log^{\#A}{x} - 
%         \int_{1}^{x}\frac{\pi_A(t)\cdot \#A \cdot \log^{\#A-1}{t}}{t}dt$$
%     and
%     $$\pi_A(x) = \frac{\upvartheta_A(x)}{\log^{\#A}{x}} + 
%         \int_{2}^{x}\frac{\#A}{t\log^{\#A+1}{t}}dt.$$
% \end{proposition}
% \begin{proof}
%     The proof follows exactly the path of Theorem 4.3 of \cite{Apostol},
%     and so is omitted.
% \end{proof}

% \begin{proposition}
%     The following relations are logically equivalent:
%     $$\pi_A(x) \sim \frac{C_Ax}{\log^{\#A}{x}}.$$
%     $$\upvartheta_A(x) \sim C_Ax.$$
% \end{proposition}

% \begin{definition}
%     $$\Psi_A(x) = \sum_{n \le x} \Lambda(n)w_A(n).$$
% \end{definition}
% \begin{theorem}
%     For any admissible k-tuple $A$, we have
%     $$\Psi_A(x) = \Psi(x) - \#S_A\log x + O(1).$$
% \end{theorem}
% \begin{proof}
%     We have
%     \begin{align*}
%     \Psi(x)-\Psi_A(x)
%     &= \sum_{n\le x} \Lambda(n)(1-w_A(n)).
%     \end{align*}
    
%     Now, for composite $n$ we have $\Lambda(n)=0$, and for $n=p^k$ where 
%     $p\not\in S_A$, we have $1-w_A(n)=0$, so the summand only takes
%     on values at powers of $p\in S_A$, and we have
%     \begin{align*}
%     \Psi(x)-\Psi_A(x)
%     &= \sum_{p\in S_A}\log(p)\sum_{k \le \log_p x} (1-w_A(p)^k) \\
%     &= \sum_{p\in S_A}\log(p)\sum_{k \le \log_p x} 1
%         + O\left(\sum_{p\in S_A}\log(p)\sum_{k \le \log_p x} w_A(p)^k\right) \\
%     &= \sum_{p\in S_A}\log(p)\sum_{k \le \log_p x} 1
%         + O(1) \\
%     &= \sum_{p\in S_A}\log(p)(\log_p(x) + O(1)) + O(1) \\
%     &= \sum_{p\in S_A}\log(x) + O(1) \\
%     &= \#S_A\log(x) + O(1).
%     \end{align*}

% \end{proof}
% \begin{proof}[Silly proof]
%     We apply Perron's formula and obtain the following where the stars
%     indicate the summation should be multiplied by 1/2
%     when $x\in\mathbb{Z}$.
%     \begin{align*}
%         \Psi^*_A(x) &= \sum_{n\le x}^*\Lambda(n)w_A(n) \\
%         &= \frac{1}{2\pi i}\int_{c-i\infty}^{c+i\infty}\Di(\Lambda\cdot w_A)(z)
%             \frac{x^z}{z}dz,
%     \end{align*}
%     where $c>1$.
%     Because $w_A$ is completely multiplicative, we have 
%     \begin{align*}
%     \Di(\Lambda\cdot w_A)(z) 
%     &= -\frac{d}{dz}\log(\Di(w_A)(z)) \\
%     &= -\frac{d}{dz}\log(\zeta(z)M_A(z)) \\
%     &= -\frac{d}{dz}\log(\zeta(z)) - \frac{d}{dz}\log(M_A(z)) \\
%     &= \Di(\Lambda)(z) - \frac{d}{dz}\log(M_A(z)) \\
%     \end{align*}

%     Substituting this into our integral, we then have
%     \begin{align*}
%         \Psi^*_A(x)
%         &= \frac{1}{2\pi i}\int_{c-i\infty}^{c+i\infty}\left[
%             \Di(\Lambda)(z)
%             - \frac{d}{dz}\log M_A(z)\right]
%             \frac{x^z}{z}dz \\
%         % &= \Psi^*(x)
%         %     - \frac{1}{2\pi i}\int_{c-i\infty}^{c+i\infty}\frac{M'_A}{M_A}(z)
%         %     \frac{x^z}{z}dz \\
%         &= \Psi^*(x)
%             - \frac{1}{2\pi i}\int_{c-i\infty}^{c+i\infty}
%                 \left[\frac{d}{dz}\log M_A(z)\right]\frac{x^z}{z} dz.
%     \end{align*}

%     The derivative can be expanded as
%     \begin{align*}
%         \frac{d}{dz}\log M_A(z) 
%         &= \frac{d}{dz} \sum_{p\in S_A}\left[\log(1-p^{-z}) - \log(1-w_A(p)p^{-z})\right] \\
%         &= \sum_{p\in S_A}\left[\frac{\log(p) p^{-z}}{1-p^{-z}} - \frac{w_A(p)\log(p) p^{-z}}{1-w_A(p)p^{-z}}\right] \\
%         &= \sum_{p\in S_A}\left[\log(p)\sum_{k\ge 1}p^{-kz}(1-w_A^k(p))\right]. \\
%     \end{align*}
%     (TODO: Is this a stupid way to eval this? There's probably a more direct route, but
%     integrating series makes everything easy except for convergence.)

%     Now, substituting this into our above integral and distributing the integral
%     over all sums (TODO convergence), we have
%     \begin{align*}
%         \Psi^*_A(x) 
%         &= \Psi^*(x)
%             - \frac{1}{2\pi i}
%                 \sum_{p\in S_A} \log p
%                 \sum_{k\ge 1}
%                 \int_{c-i\infty}^{c+i\infty} \frac{(x/p^k)^z}{z}dz.
%     \end{align*}

%     When $\log_p(x)<k$ we have $x/p^k<1$, and we may consider the integral to
%     be closed along a right semicircle (just trust me dawg, but actually this
%     is a critical detail).
%     The integrand is analytic within that half-plane, so the integrals are all 0.

%     When $\log_p(x)>k$, we may then consider the integral to be along a left
%     closed semicircle, so the integral evaluates to $2\pi i \cdot \Res((x/p^k)^z/z,0)
%     = 2\pi i$.

%     When $\log_p(x)=k$, the integral may be directly evaluated on the path
%     $z=c+it$ as the vertical bound go symmetrically to infinity, and yields $\pi i$.

%     Accounting for each case, we have
%         \begin{align*}
%         \Psi^*_A(x) 
%         &= \Psi^*(x)
%             - \sum_{p\in S_A} \log p
%               \sum^*_{k \le \log_p x} 1 \\
%         &= \Psi^*(x)
%             - \sum_{p\in S_A} \log p
%               (\log_p(x) + O(1)) \\
%         &= \Psi^*(x) - \#S_A\log x + O(1).
%     \end{align*}

%     Now, the change in order $d(x)$ created by removing the stars when $x\in\mathbb{Z}$ is
%     \begin{align*}
%     d(x) = (\Psi(x)-\Psi_A(x)) - (\Psi^*(x)-\Psi^*_A(x))
%     & = \frac{1}{2}\Lambda(x)(1-w_A(x)).
%     \end{align*}
%     If $x$ is a power of a prime $p\in S_A$, $d$ is bounded by the max of
%     the logs of those finitely many primes.
%     If $x$ is a power of any other prime, we have $w_A(x)=1$, so $d=0$.
%     If $x$ is any other number, $\Lambda(x)=0$, so $d=0$.
%     Then $d(x)=O(1)$, and the result follows.
% \end{proof}

% Combining this with the Prime Number Theorem in the form $\Psi(x) \sim x$ gives
% $\Psi_A(x) \sim x$.

\begin{enumerate}
    \item
        We have $\nu_A\ast\phi_A = N = \phi_A^0$.
        That is, $\nu_A$ jumps between levels of phi.
        What functions jump between other levels?

    \item
        Can the quota totient formulae be inverted to express higher quotas in
        terms of lower?
        Are there corresponding formulae relating unit functions for quota totients?

    \item
        Double check Granville actually implies the thing.
        Oh, just noticed he is the pretentious theory guy.

    \item 
        We should check if any results require finiteness on k-tuple.
        Also check each theorem for admissibility requirements.
\end{enumerate}

\pagebreak
\printbibliography

\end{document}

